\documentclass{beamer}
%% fink install texlive
\usetheme{Warsaw}
\usepackage{amsmath,amssymb}
\usepackage{stmaryrd}
\usepackage{graphicx}
%\usepackage{haskell}
\usepackage{code}
%\usepackage{proof}
\usepackage{theorem}
\usepackage{pstricks} 


\usepackage{tikz}
\usetikzlibrary{shapes.geometric, arrows, shapes}

\tikzstyle{startstop} = [rectangle, rounded corners, minimum width=1.5cm, minimum height=0.5cm,text centered, draw=black]
\tikzstyle{io} = [trapezium, trapezium left angle=70, trapezium right angle=110, minimum width=1.5cm, minimum height=0.5cm, text centered, draw=black]
\tikzstyle{process} = [rectangle, minimum width=1.5cm, minimum height=0.5cm, text centered, draw=black]
\tikzstyle{decision} = [diamond, minimum width=1.5cm, minimum height=0.5cm, text centered, draw=black, aspect=3]
\tikzstyle{arrow} = [thick,->,>=stealth]
\tikzstyle{doublearrow} = [thick,<->,>=stealth]
\tikzstyle{db} = [cylinder, shape border rotate=90, draw,minimum
height=1.5cm, minimum width=2cm, aspect=0.3]
\tikzstyle{www} = [cloud, cloud puffs=15.7, cloud ignores
aspect,minimum width=5cm, minimum height=2cm, align=center, draw]
\tikzstyle{blank} = [rectangle] 
\tikzstyle{container} = [rectangle, rounded corners, minimum
width=7cm, minimum height=7cm,text centered, draw=black]

\tikzstyle{darrow} = [thick,<->,>=stealth]

\usepackage{listings}


% Martin's macros
   

% groups

%%\newtheorem{lemma}{Lemma}
%%\newtheorem{definition}{Definition}
%%\newtheorem{theorem}{Theorem}
%%\newtheorem{corollary}{Corollary}
%%\newtheorem{remark}{Remark}
%%\newtheorem{conjecture}{Conjecture}
%%\newtheorem{example}{Example}

% arrays

\newcommand{\bi}{\begin{array}[t]{@{}l@{}}}
\newcommand{\ei}{\end{array}}
\newcommand{\ba}{\begin{array}}
\newcommand{\ea}{\end{array}}
\newcommand{\bda}{\[\ba}
\newcommand{\eda}{\ea\]}
\newcommand{\bp}{\begin{quote}\tt\begin{tabbing}}
\newcommand{\ep}{\end{tabbing}\end{quote}}
\newcommand{\set}[1]{\left\{
    \begin{array}{l}#1
    \end{array}
  \right\}}
\newcommand{\sset}[2]{\left\{~#1 \left|
      \begin{array}{l}#2\end{array}
    \right.     \right\}}


% symbols

\newcommand{\good}{good}
 
\newcommand{\err}{{\bf W}}
\newcommand{\VV}{{\cal V}} 
\newcommand{\GT}{{\cal T}_G}
\newcommand{\KK}{{\cal K}}
\newcommand{\TT}{{\cal T}}
\newcommand{\TC}{\mbox{\it TC}}
\newcommand{\UL}{\mbox{UL}}
\newcommand{\sUL}{\mbox{\scriptsize UL}}
\newcommand{\true}{\mbox{\sf true}}
\newcommand{\Int}{\mbox{\it Int}}
\newcommand{\Bool}{\mbox{\it Bool}}
\newcommand{\SF}{\mbox{$\cal S$}}
\newcommand{\tauvec}{\bar{\tau}}
\newcommand{\tvec}{\bar{t}}
\newcommand{\kvec}{\bar{k}}
\newcommand{\xvec}{\bar{x}}
\newcommand{\muvec}{\bar{\mu}}
\newcommand{\alphavec}{\bar{\alpha}}
\newcommand{\betavec}{\bar{\beta}}
\newcommand{\gammavec}{\bar{\gamma}}
\newcommand{\thetavec}{\bar{\theta}}
\newcommand{\deltavec}{\bar{\delta}}
\newcommand{\bvec}{\bar{b}}
\newcommand{\typo}{\Gamma}
\newcommand{\tenv}{\Gamma}
\newcommand{\typoinit}{\typo_0}
\newcommand{\static}{S}
\newcommand{\dynamic}{D}
\newcommand{\sstatic}{{\scriptstyle S}}
\newcommand{\sdynamic}{{\scriptstyle D}}
\newcommand{\BTA}{\mbox{BTA}}
\newcommand{\baseshape}{\bot}
\newcommand{\TCT}{P}
\newcommand{\Haskell}{H98}
\newcommand{\TCTinit}{P_o}

% figures, rules


\newcommand{\tlabel}[1]{\mbox{(#1)}}
\newcommand{\fig}[3]
        {\begin{figure*}[t]#3\
        \caption{\label{#1}#2}\ \hrulefill\ \end{figure*}}
\newcommand{\figurebox}[1]
        {\fbox{\begin{minipage}{\textwidth} #1 \end{minipage}}}
\newcommand{\boxfig}[3]
        {\begin{figure*}\figurebox{#3\caption{\label{#1}#2}}\end{figure*}}
\newcommand{\myirule}[2]{{\renewcommand{\arraystretch}{1.2}\ba{c} #1
                      \\ \hline #2 \ea}}

% relations, operators

\newcommand{\sat}{\mbox{\it sat}}
\newcommand{\entail}{\mbox{\it entail}}
\newcommand{\unambig}{\mbox{\it unambig}}
\newcommand{\complete}{\mbox{\it complete}}

\newcommand{\projexclusion}[1]{\mbox{$\bar{\exists}#1$}}
\newcommand{\funinst}{\mbox{\it inst}}
%%\renewcommand{\inst}{\, \vdash^{\scriptstyle i} \,}
%%\newcommand{\inst}{\, \vdash \,}
\newcommand{\ileq}{\preceq}
\newcommand{\ieq}{\simeq}
\newcommand{\topannot}[1]{|#1|}
\newcommand{\gendelta}{\mbox{\it gen}_{\delta}}
\newcommand{\genbeta}{\mbox{\it gen}_{\beta}}
\newcommand{\gen}{\mbox{\it gen}}
\newcommand{\normalize}{\mbox{\it normalize}}
\newcommand{\fail}{\mbox{\it fail}}
\newcommand{\fixpt}{\mbox{${\cal F}$}}
\newcommand{\testpa}{\mbox{${\cal T}$}}
\newcommand{\addpa}{\mbox{${\cal A}$}}
\newcommand{\matchpa}{\mbox{${\cal M}$}}
\newcommand{\refmatchpa}{\mbox{${\cal R}$}}
\newcommand{\wft}[1]{\mbox{\it wft}(#1)}
\newcommand{\tv}{\mbox{\it fv}}
\newcommand{\stv}{\mbox{\it semfv}}
\newcommand{\range}{\mbox{\it range}}
\newcommand{\lub}{\wedge}
\newcommand{\tyrel}[2]{(#1 \preceq #2)}
\newcommand{\eqty}[2]{#1=#2}
%%\newcommand{\ts}{\, \vdash \,}
\newcommand{\turns}{\, \vdash \,}
\newcommand{\tsinst}{\, \vdash^i \,}
\newcommand{\tsttv}{\, \vdash^{ttv} \,}
\newcommand{\tsUL}{\, \vdash_{\mbox{\scriptsize UL}} \,}
\newcommand{\trc}{\, \vdash_{\scriptstyle \id{inf}} \,}
\newcommand{\trcp}{\, \vdash_{\scriptstyle inf}' \,}
\newcommand{\genrel}{\, \vdash^{\scriptstyle gen} \,}
\newcommand{\asubty}[2]{(#1 \leq_a #2)}
\newcommand{\ssubty}[2]{(#1 \leq_s #2)}
\newcommand{\fsubty}[2]{(#1 \leq_f #2)}
\newcommand{\subty}[2]{(#1 \leq #2)}
\newcommand{\doubleb}[1]{[\![ #1 ]\!]}
\newcommand{\id}[1]{\mbox{\it #1}}


\newcommand{\mynote}[1]{$\spadesuit${\bf #1}$\clubsuit$}

% spacing

\newcommand{\sgap}{\quad}
\newcommand{\bgap}{\quad\quad}

% reserved words

\newcommand{\mathem}{\sf}
\newcommand{\IN}{\mbox{\mathem in}}
\newcommand{\LET}{\mbox{\mathem let}}
\newcommand{\MLET}{\mbox{\mathem mlet}}
\newcommand{\LETREC}{\mbox{\mathem letrec}}
\newcommand{\FIX}{\mbox{\mathem fix}}
\newcommand{\CASE}{\mbox{\mathem case}}
\newcommand{\TYCASE}{\mbox{\mathem typecase}}
\newcommand{\ELSE}{\mbox{\mathem else}}
\newcommand{\IF}{\mbox{\mathem if}}
\newcommand{\OF}{\mbox{\mathem of}}
\newcommand{\THEN}{\mbox{\mathem then}}
\newcommand{\INST}{\mbox{\mathem instance}}
\newcommand{\OVER}{\mbox{\mathem overload}}
\newcommand{\CLASS}{\mbox{\mathem class}}
\newcommand{\WHERE}{\mbox{\mathem where}}


%%% Local Variables: 
%%% mode: latex
%%% TeX-master: "bta"
%%% End: 


\newcommand{\rev}[1]{{#1}^{\mbox{\scriptsize r}}}
\newcommand{\lquo}{\backslash}
\newcommand{\rquo}{/}
\newcommand{\diff}{-}
\newcommand{\isect}{\cap}
\newcommand{\mymid}{~\redtxt{\mid}~}
\newcommand{\lconf}[2]{#1:{\tt #2}}

\newcommand{\ignore}[1]{}

\newcommand{\magtxt}[1]{{\magenta #1}}
\newcommand{\redtxt}[1]{{\red #1}}
\newcommand{\bluetxt}[1]{{\blue #1}}
\newcommand{\greytxt}[1]{{\gray #1}}
\newcommand{\greentxt}[1]{{\green #1}}
\newcommand{\mmleq}{\leq}
\newcommand{\pow}{\^{}}
\newcommand{\venv}{\Delta}
\newcommand{\mleq}{\mbox{\tt leq}}
\newcommand{\mas}{\mbox{\tt as}}
\newcommand{\subt}{<\!:}

\newcommand{\Nturns}{\, \vdash_{\mbox{\tiny lnf}} \,}

\newcommand{\deriv}[2]{#1\backslash #2}
\newcommand{\pderiv}[2]{pd_{#2}(#1)} %% {#1\backslash_p #2}
\newcommand{\pderivFail}[2]{pd^{err}_{#2}(#1)} %% {#1\backslash_p #2}
\newcommand{\mkEmp}[1]{{#1}_{\epsilon}}
\newcommand{\strip}[1]{#1\downarrow}

\newenvironment{ttprog}{\begin{trivlist}\item \tt
        \begin{tabbing}}{\end{tabbing}\end{trivlist}}


\newenvironment{grammar}{%
         \begin{center} \small%
         $\begin{array}{rcll}
         }{%
         \end{array}$\end{center}\ignorespaces%
         }

\newcommand{\rem}[1]{}
\newcommand{\comment}[1]{}

\newcommand{\kl}[1]{}

\bibliographystyle{plainnat}

\begin{document}
%\lstset{language=python}
\logo{
\includegraphics[height=0.8cm]{pics/nyp-logo.eps}\vspace{230pt}
\includegraphics[height=0.6cm]{pics/sit-logo.png}\vspace{230pt}
}
\title{Introduction to Scala} 
\author{
 Kenny Zhuo Ming Lu
}
\institute[Inst.]{School of Information Technology, Nanyang Polytechnic}
\date{\today} 


\frame{\titlepage} 

%\bibliographystyle{plainnat}
%-------------------------------------------------------------------
%-------------------------------------------------------------------
\begin{frame}[fragile]
\frametitle{Agenda}
\begin{itemize}
 \item Introduction
 \item Object Oriented Programming 
 \item Functional Programming 
 \item (Generalized) Algebraic Data Type and pattern matching
 \item Functional Abstraction and Higher order function
 \item Type classes and Monads
 \item 
\end{itemize}
\end{frame}

%-------------------------------------------------------------------
%-------------------------------------------------------------------
\begin{frame}[fragile]
\frametitle{Scala is ... }

A general purpose programming language that is 
\begin{itemize}
 \item Object Oriented and Functional
 \item with distributed programming builtin
 \item with a rich type system and a huge set of libraries and eco-systems
 \item Widely used in research institutes (e.g. Spark) and industries
   (e.g. Twitter, Linkedin, Coursera Walmart, ... )
\end{itemize}
\end{frame}

%-------------------------------------------------------------------
%-------------------------------------------------------------------
\begin{frame}[fragile]
\frametitle{Hello World in a script}
Let's say we have a scala file {\tt Script.scala}
\begin{code}
println("Hello world!")
\end{code}
To execute it
\begin{code}
$ scala Script.scala
Hello world!
\end{code}
%$
\end{frame}

%-------------------------------------------------------------------
%-------------------------------------------------------------------
\begin{frame}[fragile]
\frametitle{Hello World in a console application}
Compiling Scala program into an application
\begin{code}
object Main extends App 
{
    println("Hello world!")
}
\end{code}
{\tt object} is like a singleton class.  \\
To execute it
\begin{code}
$ scalac Main.scala
$ scala Main
Hello world!
\end{code}

\end{frame}

%-------------------------------------------------------------------
%-------------------------------------------------------------------
\begin{frame}[fragile]
\frametitle{Object Oriented Programming in Scala}
It is very similar to Java
\begin{code}
class Person(n:String,i:String) {
   private val name:String = n
   private val id:String   = i
   def getName():String = name
   def getId():String = id
}
\end{code}
\begin{itemize}
 \item {\tt Person} is the name of the class, {\tt
     Person(n:String,i:String)} is a constructor.
 \item {\tt private} sets the accessibility scope of the members
 \item {\tt val} introduces a value, (immutable variable)
 \item {\tt def} introduces a function defintion.
\end{itemize}
\end{frame}

%-------------------------------------------------------------------
%-------------------------------------------------------------------
\begin{frame}[fragile]
\frametitle{Object Oriented Programming in Scala}

Class Inheritence via the {\tt extends} keyword
\begin{code}
class Student(n:String, i:String, g:Double) 
                     extends Person(n,i) {
   private var gpa = g
   def getGPA() = gpa
   def setGPA(g:Double) = {
      gpa = g
   }
}
\end{code}
\begin{itemize}
 \item {\tt var} introduces a mutable variable
 \item Most of the type annotations are optional and will be inferred
   by the type system.
\end{itemize}
\end{frame}

%-------------------------------------------------------------------
%-------------------------------------------------------------------

\begin{frame}[fragile]
\frametitle{Object Oriented Programming in Scala}
\begin{code}
class Staff(n:String, i:String, sal:Double) 
                 extends Person(n,i) {
   private var salary = sal
   def getSalary() = salary
   def setSalary(sal:Double) = 
   {
      salary = sal
   }
}
\end{code}
\end{frame}


%-------------------------------------------------------------------
%-------------------------------------------------------------------

\begin{frame}[fragile]
\frametitle{Object Oriented Programming in Scala}
Traits are like Java interfaces and abstract classes
\begin{code}
trait NightOwl {
   def stayUpLate():Unit 
}

class Student(n:String, i:String, g:Double) 
      extends Person(n,i) with NightOwl {
   private var gpa = g
   def getGPA() = gpa
   def setGPA(g:Double) = {
      gpa = g
   }
   override def stayUpLate():Unit = { 
      println("woohoo") 
   }
}
\end{code}
\end{frame}



%-------------------------------------------------------------------
%-------------------------------------------------------------------


\begin{frame}[fragile]
\frametitle{Running our first Scala program}
Scala comes with an interpretor (AKA REPL)

\begin{code}
scala> :load OOP.scala
Loading OOP.scala...
defined class Person
defined trait NightOwl
defined class Student
defined class Staff

scala> val tom = new Student("Tom", "X1235", 4.0)
tom: Student = Student@601c1dfc

scala> val martin = new Staff("Martin", "T0001", 500000.0)
martin: Staff = Staff@650fbe32

scala> tom.stayUpLate
woohoo
\end{code}
\end{frame}

%-------------------------------------------------------------------
%-------------------------------------------------------------------
\begin{frame}[fragile]
\frametitle{Functional Programming in Scala}
Recall the Haskell example
\begin{code}
data Exp =  Val Int |  Plus Exp Exp 

simp :: Exp -> Exp 
simp e@(Val v) = e
simp (Plus (Val 0) e2) = e2
simp (Plus e1 e2) = Plus (simp e1) (simp e2)
\end{code}
\end{frame}
%-------------------------------------------------------------------
%-------------------------------------------------------------------
\begin{frame}[fragile]
\frametitle{Functional Programming in Scala}
Let's rewrite in GADT style
\begin{code}
data Exp where
    Val  :: Int -> Exp
    Plus :: Exp -> Exp -> Exp

simp :: Exp -> Exp 
simp e@(Val v) = e
simp (Plus (Val 0) e2) = e2
simp (Plus e1 e2) = Plus (simp e1) (simp e2)
\end{code}
\end{frame}
%-------------------------------------------------------------------
%-------------------------------------------------------------------
\begin{frame}[fragile]
\frametitle{Functional Programming in Scala}
Algebraic data types
\begin{code}
sealed trait Exp
case class Val(v:Int) extends Exp
case class Plus(e1:Exp, e2:Exp) extends Exp

def simp(e:Exp):Exp = e match {
   case Val(v) => e
   case Plus(Val(0), e2) => e2
   case Plus(e1,e2) => Plus(simp(e1), simp(e2))
}
\end{code}
\begin{itemize}
 \item {\tt sealed} requires all extensions to the trait must be
   declared in the same module.
 \item {\tt case} makes a class ``pattern-matchable''
\end{itemize}
\end{frame}

%-------------------------------------------------------------------
%-------------------------------------------------------------------
\ignore{
\begin{frame}[fragile]
\frametitle{Functional Programming in Scala}

\begin{code}
def simpFix(e:Exp):Exp = {
   val e2 = simp(e)
   if (e == e2) {
      e
   } 
   else {
      simpFix(e2)
   }
}
\end{code}
\end{frame}
}


%-------------------------------------------------------------------
%-------------------------------------------------------------------


\end{document}

